% Definitions from the mpi-sys-macs file

% Definitions from the mpi-user-macs file
% And the definitions that we need as a result
\def\MPI/{\textsf{MPI}}    % MPI macro -- see TEX Book, pg 8,204 bloody TEX!!
\def\mpi/{\textsf{MPI}}    % MPI macro
\def\MPII/{\textsf{MPI-1}} % MPI-1 macro -- can't easily do MPI1
\def\mpii/{\textsf{MPI-1}} % MPI-1 macro
\def\MPIIDOTO/{\textsf{MPI-1.0}}   % MPI-1.0 macro - NOTE it has a O (Oh)
                                % not 0 (zero)
\def\mpiidoto/{\textsf{MPI-1.0}}   % MPI-1.0 macro - NOTE it has a O (Oh)
                                % not 0 (zero)
\def\MPIIDOTI/{\textsf{MPI-1.1}}   % MPI-1.1 macro
\def\mpiidoti/{\textsf{MPI-1.1}}   % MPI-1.1 macro
\def\MPIIDOTII/{\textsf{MPI-1.2}}  % MPI-1.2 macro
\def\mpiidotii/{\textsf{MPI-1.2}}  % MPI-1.2 macro
\def\MPIIDOTIII/{\textsf{MPI-1.3}}  % MPI-1.3 macro
\def\mpiidotiii/{\textsf{MPI-1.3}}  % MPI-1.3 macro
\def\MPIII/{\textsf{MPI-2}}        % MPI-2 macro -- can't easily do MPI2
\def\mpiii/{\textsf{MPI-2}}        % MPI-2 macro
\def\MPIIIDOTO/{\textsf{MPI-2.0}}  % MPI-2.0 macro - NOTE it has a O (Oh)
                                % not 0 (zero)
\def\mpiiidoto/{\textsf{MPI-2.0}}  % MPI-2.0 macro - NOTE it has a O (Oh)
                                % not 0 (zero)
\def\MPIIIDOTI/{\textsf{MPI-2.1}}  % MPI-2.1 macro
\def\mpiiidoti/{\textsf{MPI-2.1}}  % MPI-2.1 macro
\def\MPIIIDOTII/{\textsf{MPI-2.2}} % MPI-2.2 macro
\def\mpiiidotii/{\textsf{MPI-2.2}} % MPI-2.2 macro
\def\MPIIII/{\textsf{MPI-3}}       % MPI-3 macro -- can't easily do MPI3
\def\mpiiii/{\textsf{MPI-3}}       % MPI-3 macro
\def\MPIJOD/{\textsf{MPI-JOD}}     % MPI-JOD macro
\def\mpijod/{\textsf{MPI-JOD}}     % MPI-JOD macro
\def\MPIRT/{\textsf{MPI/RT}}       % RT MPI macros.
\def\MPIIII/{\textsf{MPI-3}}       % MPI-3 macro -- can't easily do MPI3
\def\mpiiii/{\textsf{MPI-3}}       % MPI-3 macro
\def\MPIIIIDOTO/{\textsf{MPI-3.0}}  % MPI-3.0 macro
\def\mpiiiidoto/{\textsf{MPI-3.0}}  % MPI-3.0 macro
\def\MPIIIIDOTI/{\textsf{MPI-3.1}}  % MPI-3.1 macro
\def\mpiiiidoti/{\textsf{MPI-3.1}}  % MPI-3.1 macro
\def\MPIIV/{\textsf{MPI-4}}          % MPI-4 macro
\def\mpiiv/{\textsf{MPI-4}}          % MPI-4 macro
\def\MPIIVDOTO/{\textsf{MPI-4.0}}    % MPI-4.0 macro
\def\mpiivdoto/{\textsf{MPI-4.0}}    % MPI-4.0 macro
\def\MPIIVDOTI/{\textsf{MPI-4.1}}    % MPI-4.1 macro
\def\mpiivdoti/{\textsf{MPI-4.1}}    % MPI-4.1 macro
\def\MPIV/{\textsf{MPI-5}}           % MPI-5 macro
\def\mpiv/{\textsf{MPI-5}}           % MPI-5 macro
\def\MPIVDOTO/{\textsf{MPI-5.0}}     % MPI-5.0 macro
\def\mpivdoto/{\textsf{MPI-5.0}}     % MPI-5.0 macro

\def\RMA/{\textsf{RMA}}    % RMA macro -- see TEX Book, pg 8,204 bloody TEX!!

\def\OnlyForAutomaticAnnexGeneration#1{}% deleting the content were the macro is used; but preserving it for the Annex

% pagelabel is used to get a label that is hyperlinked to the correct
% page in the PDF version. Here we can just use label
\newcommand{\pagelabel}[1]{\label{#1}}

\newcommand{\func}[1]{\textsf{#1}\mpifuncindex{#1}}
\newcommand{\mpifunc}[1]{\textsf{#1}\mpifuncindex{#1}}  % for MPI functions - language independent
% for ones that don't go in index:
\newcommand{\mpiskipfunc}[1]{\textsf{#1}}               % ... same, but not in the Function Index
\newcommand{\cfunc}[1]{\textsf{#1}}
\newcommand{\ffunc}[1]{\textsf{#1}}
\newcommand{\constskip}[1]{{\textsf{\small #1}}}
% for ones from  MPI-1
\newcommand{\consti}[1]{{\textsf{\small #1}}\index{CONST:#1}} % constants/handles - language independent
\newcommand{\constiskip}[1]{{\textsf{\small #1}}}             % ... same, but not in the Constant Index
\newcommand{\constitem}[2]{\item[\mpiconst{#1}\hfill]{#2}}
\newcommand{\constitemtwo}[3]{\item[\mpiconst{#1}, \mpiconst{#2}\hfill]{#3}}
\newcommand{\constitemthree}[4]{\item[\mpiconst{#1}, \mpiconst{#2}, \mpiconst{#3}\hfill]{#4}}
% for ones that don't go in index
\newcommand{\constskipitem}[2]{\item[\constskip{#1}\hfill]{#2}}
\newcommand{\mpiarg}[1]{\textsf{#1}}                % argument in the argument list of an MPI routine, language independent
\newcommand{\mpishortarg}[1]{\textsf{#1}}                % ... same but without \gb penalty
\newcommand{\type}[1]{\textsf{#1}\index{CONST:#1}}  % datatype handles
%
\newcommand{\gtype}[1]{\textsf{#1}} % generic (language independent) type
\newcommand{\shorttype}[1]{\textsf{#1}\index{CONST:#1}}  % ... same but without \gb panelty
\newcommand{\ctype}[1]{\texttt{#1}}                 % - and corresponding C type
\newcommand{\ftype}[1]{\texttt{#1}}                 % - and corresponding Fortran type
%
% Info is for MPI_Info predefined strings.  \infokey{keyname} and
% \infoval{keyvaluename}
\newcommand{\info}[1]{{\small\sf #1}\index{CONST:#1}}
\newcommand{\infoval}[1]{{\small\sf #1}\index{CONST:#1}}
\newcommand{\infovalshort}[1]{\infoval{#1}}
\let\infokey=\infoval
\newcommand{\infoskip}[1]{{\small\sf #1}}
\newcommand{\infokeymain}[1]{\infokeyskip{#1}\IIstringmainidx[INFOCONST=KEY]{#1}}
\newcommand{\infovalmain}[1]{\infovalskip{#1}\IIstringmainidx[INFOCONST=VAL]{#1}}
\newcommand{\infokeyskip}[1]{\IIfmatstring{#1}}
\newcommand{\infovalskip}[1]{\IIfmatstring{#1}}
%
\newcommand{\error}[1]{{\small\sf #1}\index{CONST:#1}}
\newcommand{\errorskip}[1]{{\small\sf #1}}
\newcommand{\errori}[1]{{\small\sf #1}}

\newcommand{\mpiconstfunc}[1]{\IImpifuncfmt{#1}\IIconstidx{#1}\mpifuncindex{#1}}
\newcommand{\mpiconstfuncskip}[1]{\IImpifuncfmt{#1}}
\newcommand{\mpiconstfuncmain}[1]{\IImpifuncfmt{#1}\IIconstmainidx{#1}\mpifuncmainindex{#1}}
\newcommand{\mpiconstfuncindex}[1]{\IIconstidx{#1}\mpifuncindex{#1}}

\newcommand{\IN}[0]{{\small IN}}
\newcommand{\OUT}[0]{{\small OUT}}
\newcommand{\INOUT}[0]{{\small INOUT}}

\newcommand{\exindex}[1]{\index{EXAMPLES:#1}}

\def\class{$\langle$CLASS$\rangle$}

\newcommand{\uu}[1]{\underline{\hyperpage{#1}}}
\newcommand{\mpifuncmainindex}[1]{\index{#1|uu}}
\newcommand{\mpifuncindex}[1]{\index{#1}}
\newcommand{\typedefindex}[1]{\index{TYPEDEF:#1}}
\newcommand{\cdeclindex}[1]{\index{CONST:#1}}        %  for index entry of C declarations like MPI_Comm
\newcommand{\cdeclmainindex}[1]{\index{CONST:#1|uu}}

%%%%%%%%%%%%%%%%%%%%%%%%%%%%%%%%%%%%%%%%%%%%%%%%%%%%%%%%%%%%%%%%%%%%%%%%%%%%%%
%macros for language binding: mpibind, mpifnewbind, mpifbind, and fargs:
% The binding headings can be turned on and off
% with \bindingheadingtrue and \bindingbheadingfalse. This is useful
% when all bindings are for a single language, and the heading is not needed.
\newif\ifbindingheading
\bindingheadingtrue
\newcommand{\IIbindingheading}[1]{\ifbindingheading{\noindent\textbf{#1}\par\nobreak}\fi}

\def\IIbindingcname{C binding}
\def\IIbindingfname{Fortran binding}
\def\IIbindingfnamespecial{Fortran binding (the following procedure is not available with mpif.h)}
\def\IIbindingfnewname{Fortran 2008 binding}

% This if controls whether there is a heading
\newif\ifbindinglabel
\bindinglabelfalse

% Because of how these are used in the \if tests, these definitions
% need to be a single character.
% Special case: ``s'' is really a subset of f
\def\IIbindingc{c}
\def\IIbindingf{f}
\def\IIbindingfspecial{s}
\def\IIbindingfnew{F}
\def\IIbindingdef{N}
\def\IIbindingcurlang{\IIbindingdef}
\def\IIbindingissue#1#2{\ifbindinglabel\if\IIbindingcurlang#1\else
\if\IIbindingfspecial#1\gdef\IIbindingcurlang{\IIbindingf}\else
\gdef\IIbindingcurlang{#1}\fi#2\fi\fi}
%% Use the following to reset the binding language to ``NEW'' binding
%% - E.g., A new binding will generate a heading in the correct
%% language, even if the last binding was in the same language.  Use,
%% for example, when all bindings are in C - but blocks of bindings
%% are separated in the text.
\newcommand{\bindinglangreset}{\gdef\IIbindingcurlang{N}}
%% Use \bindingislisttrue to turn *off* the spacing after the binding body
%% This is needed when there are consecutive binding blocks with no
%% intervening text. This occurs in the bindings chapter for the list
%% of handle conversion routines.
\newif\ifbindingislist \bindingislistfalse
\newcommand{\IIbindingbody}[1]{{\raggedright \hangindent 7em\hangafter=1\tt
#1 \par \ifbindingislist\else\medskip\fi}}

%%%%%%%%%%%%%%%%%%%%%%%%%%%%%%%%%%%%%%%%%%%%%%%%%%%%%%%%%%%%%%%%%%%%%%%%%%%%%%
% IMPORTANT!
% Any changes to the binding macros must be coordinated with changes
% to the binding tools (in binding-tool), buildapplang, and buildindices
%%%%%%%%%%%%%%%%%%%%%%%%%%%%%%%%%%%%%%%%%%%%%%%%%%%%%%%%%%%%%%%%%%%%%%%%%%%%%%

\newcommand{\mpibind}[2]{{\IIbindingissue\IIbindingc{\IIbindingheading{\IIbindingcname}}%
\IIbindingbody{int #1#2}}}
% mpibindmain is used for a target as the main definition. Ignore for now
\newcommand{\mpibindmain}[2]{\mpibind{#1}{#2}}

\newcommand{\mpibindnotint}[3]{{\IIbindingissue\IIbindingc{\IIbindingheading{\IIbindingcname}}%
\IIbindingbody{#1 #2#3}}}
\newcommand{\mpibindnotintmain}[3]{{\IIbindingissue\IIbindingc{\IIbindingheading{\IIbindingcname}}%
\IIbindingbody{#1 #2#3}}}

%
% This version indexes the name
% \mpiemptybindidx{routine}{returnvalue}{indexed name}
\newcommand{\mpiemptybindidx}[3]{{\par\medskip
\IIbindingissue\IIbindingc{\IIbindingheading{\IIbindingcname}}%
\IIbindingbody{#2 #1\mpifuncmainindex{#3}}}}
%
% This version does not index the name, but has same argument list
%% Used only in appLang-CNames.tex, generated ONLY by MAKE-APPLANG
\newcommand{\mpiemptybindNOidx}[3]{{\IIbindingissue\IIbindingc{\IIbindingheading{\IIbindingcname}}%
\IIbindingbody{#2 #1}}}

% special macro that avoids the int in front
% should be used for routines where don't want in index
%% UNUSED
%\newcommand{\mpiskipemptybind}[2]{{\IIbindingissue\IIbindingc{\IIbindingheading{\IIbindingcname}}%
%\IIbindingbody{#2 #1}}}

% special macro for typedef int in front of SUBROUTINE in C
%
% Note that this is indexed.
\newcommand{\mpitypedefbindmain}[2]{{\IIbindingbody{typedef int #1#2;\mpiccallbackdefindex{#1#2}}}}
\newcommand{\mpitypedefbindvoidmain}[2]{{\IIbindingbody{typedef void #1#2;\mpiccallbackdefindex{#1#2}}}}
\newcommand{\mpitypedefbind}[1]{{\IIbindingbody{typedef int #1;\mpiccallbackdefindex{#1}}}}
\newcommand{\mpitypedefbindvoid}[1]{{\IIbindingbody{typedef void #1;\mpiccallbackdefindex{#1}}}}
%
% special macro for callback prototype definitions in Fortran
\newcommand{\mpifsubbindmain}[2]{{\IIbindingbody{SUBROUTINE #1#2}\mpicallbackmainindex{#1}}}
\newcommand{\mpifsubbind}[2]{{\IIbindingbody{SUBROUTINE #1#2}}\mpifcallbackdefindex{#1}}
%
% Use \tt because we format more than a single paragraph, and \texttt
% will complain
\newcommand{\mpifnewsubbind}[2]{{\tt\noindent ABSTRACT INTERFACE\par\noindent\quad\IIbindingbody{SUBROUTINE #1#2}}\mpifcallbackdefindex{#1}}

% special macro for normal SUBROUTINE and FUNCTION definitions in Fortran
\newcommand{\mpifbind}[2]{{\IIbindingissue\IIbindingf{\IIbindingheading{\IIbindingfname}}%
\IIbindingbody{#1#2}}}
% bind should have a link to the main entry; bindmain does not
\newcommand{\mpifbindmain}[2]{\mpifbind{#1}{#2}}

% Fortran functions returning a value
\newcommand{\mpifbindret}[3]{{\IIbindingissue\IIbindingf{\IIbindingheading{\IIbindingfname}}%
\IIbindingbody{#1 #2#3}}}
% the main macro does not define a new hyperref label
\newcommand{\mpifbindretmain}[3]{\mpifbindret{#1}{#2}{#3}}

% special macro for MPI_Status_f082f | f2f08, which uses the
% "Fortran binding (the following procedure is not available with mpif.h)" label
\newcommand{\mpifbindspecial}[2]{{\IIbindingissue\IIbindingfspecial{\IIbindingheading{\IIbindingfnamespecial}}%
\IIbindingbody{#1#2}}}

\newcommand{\mpifnewbind}[2]{{\IIbindingissue{\IIbindingfnew}{\IIbindingheading{\IIbindingfnewname}}%
\IIbindingbody{#1#2}}}
% special macro for text instead of a binding; used for routines that were deprecated until MPI-2.2:
\newcommand{\mpifnewnonebind}[1]{{\IIbindingbody{\textrm{#1}}}}
% Like mpibindmain
\newcommand{\mpifnewbindmain}[2]{\mpifnewbind{#1}{#2}}
\newcommand{\mpifnewbindmainref}[2]{{\IIbindingissue{\IIbindingfnew}{\IIbindingheading{\IIbindingfnewname}}%
\IIbindingbody{#1#2}}}
% special case for Fortran functions returning a value (e.g., MPI_WTIME)
\newcommand{\mpifnewbindret}[3]{{\IIbindingissue{\IIbindingfnew}{\IIbindingheading{\IIbindingfnewname}}%
\IIbindingbody{#1 #2#3}}}
\newcommand{\mpifnewbindretmain}[3]{\mpifnewbindret{#1}{#2}{#3}}


% special macro to skip fortran binding in index
%% UNUSED
%\newcommand{\mpifskipbind}[1]{{\IIbindingbody{#1}}}

\def\fargs{\\\advance\leftskip 2em\hangindent5em\hangafter=1}
\def\fnarg{\par\hangindent5em\hangafter=1}
\def\fnfinalarg{\par\nobreak\hangindent5em\hangafter=1}

%% Use \bindindent/ to end a line and start a new line with a consistent
%% indent.  This replaces the recommendation to use
%%  \\\ \ \ \   (that last \ is followed by a space)
%% when continuing a declaration line.  This is mostly used in the Fortran
%% bindings.
%% The 7em indent is the same as for the c bindings - 5+2 (2 already for
%% fargs)
\def\bindindent/{\\\hskip5em}

\newenvironment{rationale}{\begin{list}{}{}\item[]{\it Rationale.}
}{{\rm ({\it End of rationale.})} \end{list}}

\newenvironment{implementors}{\begin{list}{}{}\item[]{\it Advice
        to implementors.}
}{{\rm ({\it End of advice to implementors.})} \end{list}}

\newenvironment{users}{\begin{list}{}{}\item[]{\it Advice to users.}
}{{\rm ({\it End of advice to users.})} \end{list}}

\newcommand{\missing}[1]{}

\newcommand{\alter}[1]{}

\newcommand{\status}[1]{}

\newcommand{\commentOut}[1]{{}}

\newcommand{\startchap}[0]{}
\newcommand{\withlinenumbers}[0]{}
\def\pdfbookmark{\eatbrack}
\def\eatbrack[{\eatnum}
\def\eatnum#1{\eatclosebrack}
\def\eatclosebrack]{\eattwo}
\def\eattwo#1#2{\relax}

\newtheorem{example}{Example}[chapter]

%\usepackage{graphicx}

%% MPI-3.1
\def\namedref#1#2{#1 \ref{#2}}
\def\sectionref#1{\namedref{Section}{#1}}
\def\textTitleFont#1{\textrm{#1}}

\def\IItermIndex#1{\IImpiindex{TERM:#1}}
\def\IItermMainIndex#1{\IImpimainindex{TERM:#1}}

\def\mpiterm#1{\IItermIndex{#1}\emph{#1}}
\def\mpitermni#1{\emph{#1}}
\def\mpitermindex#1{\IItermIndex{#1}}
\def\mpitermdef#1{\textbf{#1}\IItermMainIndex{#1}}
\def\mpitermdefni#1{\textbf{#1}}
\def\mpitermdefindex#1{\IItermMainIndex{#1}}
% Sometimes a term needs a more complex index entry; e.g., a subentry
\def\mpitermsub#1#2{\IItermIndex{#2}\emph{#1}}

% This needs to be implemented differently in HTML because the current
% implementation of list in tohtml is too simple
\newenvironment{constlist}{
    \begin{tabular}{ll}
}{\end{tabular}}

\newcommand{\noconstitem}[2]{#1&#2\\}
\newcommand{\constitem}[2]{\mpiconst{#1}&#2\\}
\newcommand{\constmainitem}[2]{\constmain{#1}&#2\\}
\newcommand{\constitemtwo}[3]{\mpiconst{#1}, \mpiconst{#2}&#3\\}
\newcommand{\constitemthree}[4]{\mpiconst{#1}, \mpiconst{#2}, \mpiconst{#3}&#4\\}
\newcommand{\constitemonly}[1]{\constitem{#1}&\\}

% mpidtypeitem is like constitem, but for MPI datatypes
\newcommand{\mpidtypeitem}[2]{\mpidtype{#1}&#2\\}
\newcommand{\mpidtypeitemonly}[1]{\mpidtypeitem{#1}&\\}
\newcommand{\mpidtypemainitem}[2]{\mpidtypemain{#1}&#2\\}

\def\mpicode#1{\textsf{#1}}
\def\mpiublb#1{\textsf{#1}\IItermIndex{#1}}

\def\mpitermtitleindex#1{\IImpiindex{TERM:#1|bold}}
\def\mpitermtitleindexmainsub#1#2{\IImpiindex{TERM:#1!#2|bold}\IImpiindex{TERM:#2|bold}}
\def\mpitermtitleindexsubmain#1#2{\IImpiindex{TERM:#1 #2|bold}\IImpiindex{TERM:#2!#1|bold}}

\def\XXX/{\textsf{XXX}}
\def\nameref#1{\ref{#1}}
\def\hyperlink#1#2{#2}
\def\url#1{\href{#1}{#1}}
\newcommand{\textoutargs}[2]{%
\funcarg{\OUT}{#1}{buffer to return the string containing #2 (string)}
\funcarg{\INOUT}{{#1}\_len}{length of the string and/or buffer for \mpiarg{#1} (integer)}
}
\newcommand{\textoutdesc}[2]{%
The arguments \mpiarg{#1} and \mpiarg{{#1}\_len} are used to return #2
as described in Section~\ref{sec:mpit:strings}.%
}
\newcommand{\sectionrange}[2]{Sections~\ref{#1} to \ref{#2}}

% MPI 4 defintions
\def\IIfontsize#1{}% Ignore fontsize control
\newcommand{\IIconstfmt}[2][0]{\textsf{\IIfontsize{#1}#2}}
\newcommand{\IImpifmt}[1]{\textsf{#1}}
%% May have to replace these, as newcommand implementation may not
%% handle this case
\newcommand{\IIconstidx}[2][CONST]{\IImpiindex{#1:#2@\IIconstfmt{#2}}}
\newcommand{\IIconstmainidx}[2][CONST]{\IImpimainindex{#1:#2@\IIconstfmt{#2}}}
\newcommand{\mpiconst}[2][0]{\mpiconstshort[#1]{#2}}
\newcommand{\mpiconstmain}[2][0]{\IIconstfmt[#1]{#2}\IIconstmainidx{#2}}
\newcommand{\mpiconstmainshort}[2][0]{\IIconstfmt[#1]{#2}\IIconstmainidx{#2}}
\newcommand{\mpiconstshort}[2][0]{\IIconstfmt[#1]{#2}\IIconstidx{#2}}
\newcommand{\mpiconstskip}[2][0]{\IIconstfmt[#1]{#2}}
\newcommand{\mpiconstmainindex}[1]{{\IIconstmainidx{#1}}}
\newcommand{\constmain}[1]{\mpiconstmain{#1}}
\newcommand{\constshortmain}[1]{\IIconstfmt[0]{#1}\IIconstmainidx{#1}}

\def\mpilegend#1{\textsf{#1}}

% MPI Datatypes - e.g., MPI_DOUBLE
% Note this is a change from the previous macros - to use the same
% font size (\small) as the other constants.
\newcommand{\mpidtype}[1]{\mpidtypeshort{#1}}
\newcommand{\mpidtypeindex}[1]{\IIconstidx{#1}}
\newcommand{\mpidtypemain}[1]{\IIconstfmt{#1}\IIconstmainidx{#1}}
\newcommand{\mpidtypemainshort}[1]{\IIconstfmt{#1}\IIconstmainidx[DATATYPE]{#1}}
\newcommand{\mpidtypeshort}[1]{\IIconstfmt{#1}\IIconstidx{#1}}
\newcommand{\mpidtypeskip}[1]{\IIconstfmt{#1}}

% MPI C types - e.g., MPI_Comm, MPI_Aint or TYPE(MPI_Status)
\newcommand{\mpictype}[1]{\IIconstfmt{#1}\IIconstidx{#1}}
\newcommand{\mpictypeshort}[1]{\IIconstfmt{#1}\IIconstidx{#1}}
\newcommand{\mpictypemain}[1]{\IIconstfmt{#1}\IIconstmainidx{#1}}
\newcommand{\mpictypeindex}[1]{\IIconstidx{#1}}
\newcommand{\mpictypemainindex}[1]{\IIconstmainidx{#1}}
\newcommand{\mpictypeskip}[1]{\IIconstfmt{#1}}
\newcommand{\mpiftype}[1]{\IIconstfmt{#1}\IIconstidx{#1}}
\newcommand{\mpiftypemain}[1]{\IIconstfmt{#1}\IIconstmainidx{#1}}
% MPI Fortran Integer kind reference.  Use \mpiftypekind{OFFSET} for
% INTEGER(KIND=MPI_OFFSET_KIND); this will index the use under
% MPI_OFFSET_KIND instead of under INTEGER(...)
% Note that this is new for MPI 4.0, and will need to be updated in
% the document.
\newcommand{\mpiftypekindshort}[1]{\IIconstfmt{INTEGER(KIND=MPI\_#1\_KIND)}%
\IIconstidx{MPI\_#1\_KIND}}
\newcommand{\mpiftypekind}[1]{\mpiftypekindshort{#1}}
\newcommand{\mpiftypekindmain}[1]{\code{INTEGER(KIND=MPI\_#1\_KIND)}%
\IIconstmainidx{MPI\_#1\_KIND}}

\newcommand{\typemain}[1]{\mpidtypemain{#1}}

\newcommand{\mpiargshort}[1]{\IImpifuncfmt{#1}}
\newcommand{\errormain}[1]{\errorskip{#1}\IIconstmainidx{#1}}
\newcommand{\errorweak}[1]{\errorskip{#1}}

\newcommand{\wpm}{World Model}
\newcommand{\spm}{Sessions Model}

\def\MPIisinitialized/{\MPI/ is \mpifuncindex{MPI\_INIT}\mpifuncindex{MPI\_INIT\_THREAD}initialized}
\def\MPIisfinalized/{\MPI/ is \mpifuncindex{MPI\_FINALIZE}finalized}
% Use these for parameters to C or Fortran functions respectively,
% esspecially MPI C and Fortran routine bindings.
\newcommand{\carg}[1]{\code{#1}}
\newcommand{\farg}[1]{\code{#1}}

%% Starter scripts like mpirun and mpiexec
\newcommand{\mpilaunch}[1]{\code{#1}\mpitermindex{#1}}
\newcommand{\mpilaunchskip}[1]{\code{#1}}
\newcommand{\mpilaunchmain}[1]{\code{#1}\IItermMainIndex{#1}}

% Miscellaneous strings, including (when fixed in the document)
% external32 and native for data representations.
\newcommand{\mpistring}[1]{\mpistringshort{#1}}
\newcommand{\mpistringshort}[1]{\IIfmatstring{#1}\IIstringidx[STRINGCONST]{#1}}
\newcommand{\mpistringskip}[1]{\IIfmatstring{#1}}
\newcommand{\mpistringmain}[1]{\IIfmatstring{#1}\IIstringmainidx[STRINGCONST]{#1}}

\def\dq{{\small\tt\char`"}}
\newcommand{\IIfmatstring}[1]{\IIconstfmt{\dq{}#1\dq}}

\newcommand{\MPIbindindex}[1]{\IImpimainindex{BCFUNCTION:#1}}
\newcommand{\MPIBCfunc}[1]{\IImpifuncfmt{#1}\IImpiindex{BCFUNCTION:#1}}
\newcommand{\MPIbindindexuse}[1]{\IImpiindex{BCFUNCTION:#1}}

% This is a temporary fix to make add some space between deprecated items
% for use in the chapter on deprecated features.
\def\deprecatedsep{\bigskip}

\newcommand{\mpicallback}[1]{\mpicallbackskip{#1}\mpicallbackindex{#1}}
\newcommand{\mpicallbackmain}[1]{\mpicallbackskip{#1}\mpicallbackmainindex{#1}}
\newcommand{\mpicallbackskip}[1]{\textsf{#1}}
\newcommand{\mpicallbackindex}[1]{\IImpiindex{TYPEDEF:#1@\textsf{#1}}}
\newcommand{\mpicallbackmainindex}[1]{\IImpimainindex{TYPEDEF:#1@\textsf{#1}}}

%% C definition (typedef) for callbacks
% As a function name, uses the default size, not \small, for the text.
% These are the MAIN (defining) index entries for these callbacks
\newcommand{\mpiccallbackdefindex}[1]{\IImpimainindexA{TYPEDEF:#1@\textsf{#1}}}
%% Fortran definition for callbacks
\newcommand{\mpifcallbackdefindex}[1]{\IImpimainindexA{TYPEDEF:#1@\textsf{#1}}}

%% References to language-specific callbacks
\newcommand{\mpiccallback}[1]{\code{#1}\mpicallbackindex{#1}}
\newcommand{\mpifcallback}[1]{\code{#1}\mpicallbackindex{#1}}
\newcommand{\mpifcallbackskip}[1]{\code{#1}}

\newcommand{\typeindex}[1]{\mpidtypeindex{#1}}
\newcommand{\ctypeshort}[1]{\texttt{#1}}
\newcommand{\constsmall}[1]{\mpiconst[1]{#1}}

\def\IImpiindex#1{\index{#1}}
\def\IImpimainindex#1{\index{#1|uu}}
% This version indexes the item in the main index UNLESS it is in the annex
% (e.g., the binding summaries)
\def\IImpimainindexA#1{\ifmpiInAnnex\IImpiindex{#1}\else\IImpimainindex{#1}\fi}

\newcommand{\ftypeshort}[1]{\texttt{#1}}
\newcommand{\mpishortfunc}[1]{\IImpifuncfmt{#1}}
\newcommand{\IIstringidx}[2][CONST]{\IImpiindex{#1:#2@\IIfmatstring{#2}}}
\newcommand{\IIstringmainidx}[2][CONST]{\IImpimainindex{#1:#2@\IIfmatstring{#2}}}
% This sets a location as a discrectionary hyphen. For HTML, ignore
\def\hyphen/{}
\newcommand{\mpifuncmain}[1]{\IImpifuncfmt{#1}\IImpimainindex{#1}}
\newcommand{\mpifuncshort}[1]{\IImpifuncfmt{#1}\IImpiindex{#1}}
\newcommand{\IImpifuncfmt}[1]{\textsf{#1}}
\newcommand{\mpifuncskip}[1]{\IImpifuncfmt{#1}}

%\newcommand{\mpifbindspecial}[1]{{\IIbindingissue\IIbindingfspecial{\IIbindingheading{\IIbindingfnamespecial}}%
%\IIbindingbody{#1}}}
\newcommand{\mpifbindspecial}[1]{{\tt #1\par\vspace{0.1in}}}

\newcommand{\paraheading}[1]{\vspace{3.0ex}\noindent\textbf{#1}}

\newcommand{\textoutoptinfo}[2]{
\MPI/ can optionally return an info object containing the default hints set
for #2.  If the
argument to \mpiarg{#1} provided by the user is the \code{NULL} pointer,
this argument is ignored, otherwise an \MPI/
implementation is required to return all hints that are supported by the
implementation for #2 and have default values specified; any user-supplied
hints that were not ignored by the implementation; and any additional hints
that were set by the implementation. If no such hints exist, a handle to a
newly created info object is returned that contains no key/value pair. The
user is responsible for freeing \mpiarg{#1} via \mpifunc{MPI\_INFO\_FREE}.
}

\newcommand{\constshort}[1]{\mpiconstshort{#1}}

\usepackage{listings}

\def\fixupbackrefformatting{}
\def\lstsmall{\relax}
